%  LaTeX support: latex@mdpi.com 
%  For support, please attach all files needed for compiling as well as the log file, and specify your operating system, LaTeX version, and LaTeX editor.

%=================================================================
\documentclass[drones,article,submit,moreauthors]{Definitions/mdpi} 
%\documentclass[preprints,article,submit,moreauthors]{Definitions/mdpi} 
% For posting an early version of this manuscript as a preprint, you may use "preprints" as the journal. Changing "submit" to "accept" before posting will remove line numbers.

% Below journals will use APA reference format:
% admsci, aieduc, behavsci, businesses, econometrics, economies, education, ejihpe, famsci, games, humans, ijcs, ijfs, investment, fclam, jintelligence, journalmedia, jrfm, jsam, languages, opermanag, peacestud, pmh, psycholint, publications, tourismhosp, youth

% Below journals will use Chicago reference format:
% arts, genealogy, histories, humanities, laws, literature, religions, risks, socsci

%--------------------
% Class Options:
%--------------------
%----------
% journal
%----------
% Choose between the following MDPI journals:
% accountaudit, acn, acoustics, actuators, addictions, addictprev, adhesives, adolescents, aerobiology, aerospace, agriculture, agriengineering, agrochemicals, agronomy, ai, aichem, aieng, aimater, aimed, aipa, air, aisens, aisoc, algorithms, allergies, alloys, amh, analog, analytica, analytics, anatomia, anesthres, animals, antibiotics, antibodies, antioxidants, applbiosci, appliedchem, appliedmath, appliedphys, applmech, applmicrobiol, applnano, applsci, aps, aquacj, archaeolstud, architecture, arm, arthropoda, asc, asi, astronautics, astronomy, atmosphere, atoms, audiolres, automation, axioms, bacteria, batteries, bdcc, beverages, biochem, bioengineering, biologics, biology, biomass, biomechanics, biomed, biomedicines, biomedinformatics, biomimetics, biomolecules, biophysica, bioresourbioprod, biosensors, biosphere, biotech, birds, blockchains, bloods, blsf, brainsci, breath, buildings, cancers, carbon, cardio, cardiogenetics, cardiovascmed, catalysts, cells, ceramics, challenges, chemengineering, chemistry, chemosensors, chemproc, children, chips, cimb, civileng, cleantechnol, climate, clinbioenerg, clinpract, clockssleep, cmd, cmtr, coasts, coatings, colloids, colorants, commodities, complexities, complications, compounds, computation, computers, condensedmatter, conservation, constrmater, cosmetics, covid, crops, cryo, cryptography, crystals, csmf, ctn, culture, curroncol, cyber, dairy, data, ddc, dentistry, dermato, dermatopathology, designs, devices, dhi, diabetology, diagnostics, dietetics, digital, disabilities, diseases, diversity, dna, drones, drugreact, drylands, dynamics, earth, ebj, ecm, ecologies, edm, eesp, electricity, electrochem, electronicmat, electronics, embodiedintell, encyclopedia, endocrines, energies, eng, engproc, entomology, entropy, environments, environremediat, epidemiologia, epigenomes, esa, est, fermentation, fibers, fintech, fire, fishes, fluids, foodeng, foods, forecasting, forensicsci, forests, fossstud, foundations, fractalfract, freshwater, fuels, future, futureinternet, futurepharmacol, futurephys, futuretransp, galaxies, gases, gastroent, gastrointestdisord, gastronomy, gels, genes, geochem, geographies, geohazards, geomatics, geometry, geosciences, geotechnics, geriatrics, germs, glacies, grainsci, grasses, green, greenhealth, gucdd, gynecolcancer, hardware, hazardousmatters, healthcare, healthcmater, hearts, hemato, hematolrep, hep, heritage, higheredu, highthroughput, horticulturae, hospitals, hydrobiology, hydrogen, hydrology, hydrometeorology, hydropower, hygiene, idr, iic, ijem, ijerph, ijgi, ijmd, ijms, ijns, ijom, ijp, ijpb, ijt, ijtm, ijtpp, ijtst, ime, immuno, industries, inflammj, informatics, information, infrastructures, inorganics, insects, instruments, inventions, iot, j, jaestheticmed, jal, japma, jcdd, jcm, jcp, jcrm, jcs, jcto, jdad, jdb, jdream, jemr, jeta, jfb, jfmk, jfth, jgbg, jgg, jimaging, jlpea, jmahp, jmmp, jmms, jmp, jmse, jne, jnt, jof, joi, joitmc, joma, jop, joptm, jor, jox, jpbi, jphytomed, jpm, jsan, jtaer, jvd, jzbg, kidneydial, kinasesphosphatases, knowledge, labmed, laboratories, lam, land, lasers, legumes, life, lights, limnolrev, lipidology, liquids, livers, logics, logistics, lubricants, lymphatics, machines, macromol, magnetism, magnetochemistry, make, marinedrugs, maritime, materials, materproc, mathematics, mca, mcb, measurements, medicina, medicines, medsci, membranes, merits, metabolites, metals, meteorology, methane, metrics, metrology, micro, microarrays, microbiolres, microelectronics, micromachines, microorganisms, microplastics, microwave, minerals, mining, mmphys, modelling, molbank, molecules, mollusca, mps, msf, mti, multimedia, muscles, nanoenergyadv, nanomanufacturing, nanomaterials, ncrna, ndt, network, neuroglia, neuroimaging, neurolint, neurosci, nitrogen, notspecified, nursrep, nutraceuticals, nutrients, obesities, occuphealth, oceans, ohbm, onco, optics, oral, organics, organoids, osteology, oxygen, pandemics, parasites, parasitologia, particles, pathogens, pathophysiology, pediatrrep, pets, pharmaceuticals, pharmaceutics, pharmacoepidemiology, pharmacy, phc, philosophies, photochem, photonics, photovoltaics, phycology, physchem, physics, physiologia, plants, plasma, platforms, pollutants, polymers, polysaccharides, populations, poultry, powders, precision, precisoncol, preprints, proceedings, processes, prosthesis, proteomes, psf, psychiatryint, psychoactives, purification, quantumrep, quaternary, qubs, radiation, rareearths, rdt, reactions, realestate, receptors, recycling, regeneration, remotesensing, reports, reprodmed, resources, rheumato, rjpm, robotics, rsee, ruminants, safety, sanpp, sca, sci, scipharm, sclerosis, seeds, sensors, separations, sexes, shi, signals, sinusitis, siuj, skins, smartcities, smartfish, sna, societies, software, soilsystems, solar, solids, spectroscj, sports, standards, stats, std, stratsediment, stresses, strokerep, superintelligence, surfaces, surgeries, suschem, sustainability, sustainfoods, sustainpackag, symmetry, synbio, systems, tae, targets, taxonomy, technologies, telecom, test, textiles, thalassrep, therapeutics, thermo, timespace, tomography, toxics, toxins, tph, transplantology, transportation, traumacare, traumas, tri, tropicalmed, universe, urbansci, uro, vaccines, vehicles, venereology, vetsci, vibration, virtualworlds, viruses, vision, waste, water, welding, wem, wevj, wild, wind, women, world, zoonoticdis

%---------
% article
%---------
% The default type of manuscript is "article", but can be replaced by: 
% abstract, addendum, article, benchmark, book, bookreview, briefcommunication, briefreport, casereport, changes, clinicopathologicalchallenge, comment, commentary, communication, conceptpaper, conferenceproceedings, correction, conferencereport, creative, datadescriptor, discussion, entry, expressionofconcern, extendedabstract, editorial, essay, erratum, fieldguide, hypothesis, interestingimages, letter, meetingreport, monograph, newbookreceived, obituary, opinion, proceedingpaper, projectreport, reply, retraction, review, perspective, protocol, shortnote, studyprotocol, supfile, systematicreview, technicalnote, viewpoint, guidelines, registeredreport, tutorial,  giantsinurology, urologyaroundtheworld
% supfile = supplementary materials

%----------
% submit
%----------
% The class option "submit" will be changed to "accept" by the Editorial Office when the paper is accepted. This will only make changes to the frontpage (e.g., the logo of the journal will get visible), the headings, and the copyright information. Also, line numbering will be removed. Journal info and pagination for accepted papers will also be assigned by the Editorial Office.

%------------------
% moreauthors
%------------------
% If there is only one author the class option oneauthor should be used. Otherwise use the class option moreauthors.

%=================================================================
% MDPI internal commands - do not modify
\firstpage{1} 
\makeatletter 
\setcounter{page}{\@firstpage} 
\makeatother
\pubvolume{1}
\issuenum{1}
\articlenumber{0}
\pubyear{2026}
\copyrightyear{2026}
%\externaleditor{Firstname Lastname} % More than 1 editor, please add `` and '' before the last editor name
\datereceived{ } 
\daterevised{ } % Comment out if no revised date
\dateaccepted{ } 
\datepublished{ } 
%\datecorrected{} % For corrected papers: "Corrected: XXX" date in the original paper.
%\dateretracted{} % For retracted papers: "Retracted: XXX" date in the original paper.
%\doinum{} % Used for some special journals, like molbank
%\pdfoutput=1 % Uncommented for upload to arXiv.org
%\CorrStatement{yes}  % For updates
%\longauthorlist{yes} % For many authors that exceed the left citation part
%\IsAssociation{yes} % For association journals

%=================================================================

\usepackage{listings}


\definecolor{codegreen}{rgb}{0,0.6,0}
\definecolor{codegray}{rgb}{0.5,0.5,0.5}
\definecolor{codepurple}{rgb}{0.58,0,0.82}
\definecolor{backcolour}{rgb}{0.95,0.95,0.92}
\lstdefinestyle{mystyle}{
    backgroundcolor=\color{backcolour},   
    commentstyle=\color{codegreen},
    keywordstyle=\color{magenta},
    numberstyle=\tiny\color{codegray},
    stringstyle=\color{codepurple},
    literate=
    {\{}{{{\color{red}\{}}}1
    {\}}{{{\color{red}\}}}}1
    {[}{{{\color{red}[}}}1
    {]}{{{\color{red}]}}}1
    {)}{{{\color{red})}}}1
    {(}{{{\color{red}(}}}1,
    basicstyle=\ttfamily\footnotesize,
    breakatwhitespace=false,         
    %breaklines=true,                 
    captionpos=b,                    
    keepspaces=true,                 
    numbers=left,                    
    numbersep=5pt,                  
    showspaces=false,                
    showstringspaces=false,
    showtabs=false,                  
    tabsize=2
}

\lstset{style=mystyle}
% Add packages and commands here. The following packages are loaded in our class file: fontenc, inputenc, calc, indentfirst, fancyhdr, graphicx, epstopdf, lastpage, ifthen, float, amsmath, amssymb, lineno, setspace, enumitem, mathpazo, booktabs, titlesec, etoolbox, tabto, xcolor, colortbl, soul, multirow, microtype, tikz, totcount, changepage, attrib, upgreek, array, tabularx, pbox, ragged2e, tocloft, marginnote, marginfix, enotez, amsthm, natbib, hyperref, cleveref, scrextend, url, geometry, newfloat, caption, draftwatermark, seqsplit
% cleveref: load \crefname definitions after \begin{document}

%=================================================================
% Please use the following mathematics environments: Theorem, Lemma, Corollary, Proposition, Characterization, Property, Problem, Example, ExamplesandDefinitions, Hypothesis, Remark, Definition, Notation, Assumption
%% For proofs, please use the proof environment (the amsthm package is loaded by the MDPI class).

%=================================================================
% Full title of the paper (Capitalized)
\Title{Regulatory Compliant Obstacle Avoidance Navigation Algorithm for Drone-Based Power Line Inspection}

% Authors, for the paper (add full first names)
\Author{ Pablo De Porcellinis $^{1}$, Xabier Olaz $^{2}$, Jes\'us Villadangos Alonso $^{2}$ and Daniel Al\'aez $^{2,}$*}

% MDPI internal command: Authors, for metadata in PDF
\AuthorNames{Pablo De Porcellinis, Xabier Olaz, Jes\'us Villadangos Alonso and Daniel Al\'aez}

% Affiliations / Addresses
\address{%
$^{1}$ \quad FuVeX Civil SL and Public University of Navarre (UPNA), Spain; p.porcellinis@fuvex.com\\
$^{2}$ \quad Department of Mathematical Engineering and Computer Science, Public University of Navarre (UPNA), Spain}

% Contact information of the corresponding author
\corres{Correspondence: daniel.alaez@unavarra.es}

% Current address and/or shared authorship
%\firstnote{Current address: Affiliation.}  
% Current address should not be the same as any items in the Affiliation section.

%\secondnote{These authors contributed equally to this work.}
% The commands \thirdnote{} till \eighthnote{} are available for further notes.

%\simplesumm{} % Simple summary

%\conference{} % An extended version of a conference paper

% Abstract (Do not insert blank lines, i.e. \\) 
\abstract{Most existing UAV obstacle-avoidance algorithms focus on collision-free flight, geometric efficiency, or dynamic feasibility. Much less work addresses the European regulatory logic that an autonomous route should avoid exposing uninvolved people and other protected ground parties to unnecessary overflight risk. This paper presents Path Planning Obstacle Avoidance Regulatory Compliant Evasion, a deterministic real-time local replanning method for waypoint-preserving obstacle avoidance in autonomous power-line inspection missions. The method is implemented and evaluated inside a simulated real environment in Unreal, where it operates together with a waypoint-driven flight controller, a YOLO-based perception front-end, and a  monocular obstacle geopositioning. In the current implementation, civilian-like detections and nearby entities are converted into protected local no-fly regions that trigger conservative detours instead of nominal overflight. The proposed method combines the EASA-motivated no-overflight objective with deterministic online local replanning and autonomous power-line inspection integration. In the validated static-obstacle campaign, the system remains dormant while no obstacle crosses the speed-adaptive reaction horizon, computes a bounded local bypass over an $81 \\times 81$ occupancy grid with a hard 12 m safety radius when a conflict is confirmed, and returns to the original corridor once the protected region is cleared, adding only a short conservative detour to the nominal mission. The presented results show the proposed solution as a reproducible, low-overhead, and auditable local replanning method that operationalizes part of the European ground-risk logic at the online controller level, while the simulated workflow provides the integration and validation environment used to characterize its behavior.}

% Keywords
\keyword{Ground risk; local replanning; Regulation Compliance;  safety-critical navigation; SORA; UAV inspection; uninvolved people} 

% The fields PACS, MSC, and JEL may be left empty or commented out if not applicable
%\PACS{J0101}
%\MSC{}
%\JEL{}

%%%%%%%%%%%%%%%%%%%%%%%%%%%%%%%%%%%%%%%%%%
% Only for the journal Diversity
%\LSID{\url{http://}}

%%%%%%%%%%%%%%%%%%%%%%%%%%%%%%%%%%%%%%%%%%
% Only for the journal Applied Sciences
%\featuredapplication{Authors are encouraged to provide a concise description of the specific application or a potential application of the work. This section is not mandatory.}
%%%%%%%%%%%%%%%%%%%%%%%%%%%%%%%%%%%%%%%%%%

%%%%%%%%%%%%%%%%%%%%%%%%%%%%%%%%%%%%%%%%%%
% Only for the journal Data
%\dataset{DOI number or link to the deposited data set if the data set is published separately. If the data set shall be published as a supplement to this paper, this field will be filled by the journal editors. In this case, please submit the data set as a supplement.}
%\datasetlicense{License under which the data set is made available (CC0, CC-BY, CC-BY-SA, CC-BY-NC, etc.)}

%%%%%%%%%%%%%%%%%%%%%%%%%%%%%%%%%%%%%%%%%%
% Only for the journal BioTech, Fishes, Neuroimaging and Toxins
%\keycontribution{The breakthroughs or highlights of the manuscript. Authors can write one or two sentences to describe the most important part of the paper.}

%%%%%%%%%%%%%%%%%%%%%%%%%%%%%%%%%%%%%%%%%%
% Only for the journal Encyclopedia
%\encyclopediadef{For entry manuscripts only: please provide a brief overview of the entry title instead of an abstract.}

%%%%%%%%%%%%%%%%%%%%%%%%%%%%%%%%%%%%%%%%%%
% Different journals have different requirements. Please check the specific journal guidelines in the "Instructions for Authors" on the journal's official website.
%\addhighlights{yes}
%\renewcommand{\addhighlights}{%
%
%\noindent This is an optional section, the goal of which is to increase the discoverability and readability of the article via search engines and other scholars. Highlights should not be a copy of the abstract, but a simple overview allowing the reader to quickly and simply find out what the article is about and what can be cited from it. Each of these parts should be up to 2~bullet points.\vspace{3pt}\\
%\textbf{What are the main findings?}
% \begin{itemize}[labelsep=2.5mm,topsep=-3pt]
% \item First bullet.
% \item Second bullet.
% \end{itemize}\vspace{3pt}
%\textbf{What are the implications of the main findings?}
% \begin{itemize}[labelsep=2.5mm,topsep=-3pt]
% \item First bullet.
% \item Second bullet.
% \end{itemize}
%}

%%%%%%%%%%%%%%%%%%%%%%%%%%%%%%%%%%%%%%%%%%
\begin{document}


\section{Introduction}

Unmanned Aerial Vehicles (UAVs) are becoming a practical tool for infrastructure inspection, surveillance, and monitoring, with a growing impact in agriculture, critical infrastructure, and other industrial environments \cite{imran2025environmental,marut2024surveillance}. Power-line inspection is one of the most relevant use cases because it combines long-range navigation, repeated geometric structures, tight mission corridors, and the need to react to unexpected obstacles without compromising the inspection objective. The operational experience of companies such as FuVeX in this scenario helps to identify the remaining gaps and motivates the solution presented in this paper. 

The planning difficulty is twofold. On one hand, the aircraft must preserve a global mission structure while reacting online to conflicts that appear only at close range; purely global planners handle this well only when the environment is static and fully known, and degrade quickly under late obstacle appearance or perception uncertainty \cite{debnath2024review,meng2025advances}. On the other hand, fully learned end-to-end policies, although reactive, are difficult to audit in safety-critical aeronautical contexts where traceability and bounded behavior are essential design properties \cite{EASA_AI_Roadmap_2.0}.

European regulation adds a further layer of difficulty. The operations of civil unmanned aircraft in the EU are governed by Implementing Regulation (EU) 2019/947 and Delegated Regulation (EU) 2019/945, and the acceptable means of compliance now embed the JARUS SORA 2.5 methodology for the `specific' category \cite{EU2019_947,EU2019_945,JARUS_SORA_2_5_2024,EASA_AMC_GM_SORA25_2025}. A direct consequence is that the goodness of a trajectory cannot be judged on geometry alone: a route that is perfectly flyable in kinematic terms may still be operationally unacceptable if it needlessly overflies uninvolved persons or sensitive areas on the ground. The planner must therefore answer not only whether the aircraft \emph{can} pass over a point, but whether it \emph{should} in a risk-aware sense. Several authors have begun to capture this concern using urban risk maps, third-party risk models, SORA-informed cost functions, and target-level-of-safety constraints \cite{primatesta2019riskaware,primatesta2020groundrisk,pang2022integrated,tang2023thirdparty,zhang2023groundrisk,heinze2023sora,brassel2025riskaware}. Yet these contributions are, for the most part, strategic or pre-flight route designers operating at city or corridor scale, rather than reactive local replanners acting inside an inspection mission that is already underway.

 The present paper does not encode the full legal tables of open-category separation or the complete SORA mitigation process within the planner. It implements only the controller-level step that these regulatory mechanisms motivate: live civilian-like detections are converted into temporary protected regions that the UAV should avoid crossing.


This paper focuses on Path Planning and 
Obstacle Avoidance Regulatory Compliant Evasion (PORCE) as the main scientific contribution. %The method is a local conflict-resolution system designed for the validated \texttt{SIMULATION} workflow available in the Deep-AeroTwin repository \cite{Deep_Aero_Twin}, where a mission file, a flight controller, an onboard-like perception stack, and a local planner operate together in closed loop. 
Its role is simple: whenever an obstacle enters a conservative reaction horizon, the method temporarily inhibits nominal navigation, computes a bounded local path toward the current mission waypoint, and returns control to the original route as soon as the local conflict disappears. The contribution of the paper is therefore not obstacle avoidance itself, but the integration of a regulatory-aware no-overflight logic for uninvolved third parties inside a deterministic online local replanner. We present the method as a state-of-the-art contribution in this specific EASA-relevant problem, and we do not claim to supersede the broader literature on generic UAV avoidance or strategic risk-aware routing.

The main contributions of the paper are the following.
\begin{itemize}
    \item The definition of PORCE as a deterministic, waypoint-preserving local replanning method for real-time UAV obstacle avoidance in power-line inspection missions, explicitly motivated by European ground-risk and no-overflight concerns.
    \item A formal and reproducible description of how the method is integrated with perception, mission execution, and evaluation of results inside a validated simulation environment, together with an end-to-end validation against \emph{static} obstacles (power-line assets and livestock). The treatment of \emph{moving} obstacles, with statistical and audited evidence, is reported in a companion study.
    
    %\item An evidence-driven evaluation of the method grounded on repository artifacts, including controlled ablation logs and an audited case study with figures, tables, and derived metrics.

\end{itemize}

%\section{State of the Art, Regulatory Concerns and Research Gap}

\subsection{State of the Art}
Vision-based UAV navigation has been widely studied as a perception-driven alternative to map-only planning, with applications ranging from localization to obstacle detection and local trajectory generation. Although these techniques are often motivated by scenarios where Global Navigation Satellite Systems (GNSS) signals are unavailable, unreliable, or degraded, the literature has demonstrated that vision-based approaches are not limited to such conditions \cite{arafat2023vision,huang2025investigation}. However, there are general reviews on UAV navigation that integrate this type of techniques and consistently distinguish between global and local path planning paradigms. Global path planning methods assume a sufficiently known environment and compute a static solution in advance based on prior maps or predefined constraints. Classical search algorithms such as Dijkstra or A* have been employed within this paradigm and can perform effectively in fully mapped and static environments \cite{lavalle2006planning}.However, their use as stand-alone solutions without integration with complementary perceptual or reactive components presents significant limitations in real-world missions such as power-line inspection or infrastructure monitoring. In contrast, local path planning approaches generate or adapt trajectories online while the vehicle is moving, allowing immediate reactions to newly perceived environmental elements encountered during flight \cite{debnath2024review,meng2025advances}. More specifically, recent surveys in 2025 and 2026 confirm that the dominant design goals remain collision avoidance, dynamic feasibility, and computational efficiency, while explicit treatment of ground exposure of third-parties is still comparatively rare in runtime planning \cite{merei2025survey,nafees2026review}. The proposed method is explicitly positioned in the online local-replanning class.


%However, there is a second distinction that matters for this paper and that was often implicit in earlier project notes: much of the academic literature is still evaluated under what may be called a \emph{toy-problem bias}. In those formulations, the navigation problem is simplified enough to validate a planner mathematically, but the same simplifications weaken direct transfer to operational infrastructure inspection. Three simplifications are especially consequential here.
%\begin{itemize}[leftmargin=1.3em]
%    \item \textbf{Obstacle homogeneity.} Obstacles are often represented as interchangeable geometric primitives or generic occupied cells. In inspection practice, this is not operationally neutral: a tower, a cow, and an uninvolved person may all occupy space, but they do not carry the same legal or operational meaning.
%    \item \textbf{Binary safety.} Many planners implicitly treat a mission as successful as long as physical collision is avoided. For European civil UAS operations, that criterion is too weak because a near overflight of uninvolved people may already be operationally unacceptable even when no impact occurs \cite{EU2019_947,JARUS_SORA_2_5_2024,EASA_AMC_GM_SORA25_2025}.
%    \item \textbf{Free-space permissiveness.} Conventional local planning usually assumes that any geometrically free air volume is candidate flight space. In European operations, parts of that free volume may still be undesirable or prohibited once uninvolved people, sensitive ground areas, or risk buffers are considered.
%\end{itemize}
%This is why the present paper does not frame the problem as generic collision avoidance alone. It frames it as online mission execution under geometric and regulatory pressure at the same time.


For power-line inspection, the local planning problem is more operational than academic. The vehicle must continue progressing through a mission corridor, while structures, animals, vehicles, or uninvolved people may appear in the neighborhood of the route. Power-line inspection has therefore developed its own system literature, from UAV-image approaches for corridor measurement and obstacle localization \cite{zhang2017powerline} to line detection in cluttered scenes \cite{zhang2019powerlines}, LiDAR-supported automatic inspection \cite{guan2021uavlidar}, dual-view stereovision for reconstruction and clearance estimation \cite{zhou2022dualview}, and dedicated reviews of UAV-enabled power-line inspection \cite{mendu2025state}. On the perception side, Bellou \emph{et al.} report aerial inspection based on YOLOv8 for towers, conductors and insulators \cite{bellou2024aerial}, while Faisal \emph{et al.} synthesize more than 70 studies and identify persistent gaps in multimodal fusion and deployment realism \cite{faisal2025deep}. On the planning side, Huang \emph{et al.} propose a hybrid RRT*+DWA method for dynamic obstacle avoidance in UAV-based power inspection \cite{huang2025hybrid}. These works are highly relevant, but their central objective is still efficient and safe inspection autonomy, not explicit no-overflight treatment of uninvolved people under European operational logic.

\subsection{Regulatory Concerns: Uninvolved People and Ground Risk}

To clarify the motivation of the proposed solution, this section briefly summarizes the UAS regulatory framework defined by Regulations (EU) 2019/947 and 2019/945 and applied through AMC/GM and SORA 2.5 \cite{EU2019_947,EU2019_945,JARUS_SORA_2_5_2024,EASA_AMC_GM_SORA25_2025}. The key concept around regulation is operational risk, so there are topics within this framework that transform navigation into an optimization problem with hard restrictions. The concept of an uninvolved person plays a significant role in the segmentation of categories. EASA defines it as a person who does not participate in the UAS operation or who is not aware of the instructions and safety precautions given to the UAS operator. In contrast to a power-line tower or cable, an "uninvolved person" generates a risk field in the surrounding space that adds new constraints inside the route planning, so that route geometry cannot be judged independently from who may be under the aircraft and under what conditions. 

As the intrinsic risk of an operation grows, the treatment of uninvolved people becomes a central regulatory concern, and it is managed through the Ground Risk Buffer: a spatial margin surrounding the Operational Volume that is controlled through a combination of operational constraints, technical mitigations, and procedural measures. From the planner's perspective, this buffer is part of the safety envelope that the trajectory must respect. Its extent is not fixed: according to the SORA Annex~F formulation \cite{JARUS_SORA_2_5_2024}, it grows with the kinetic energy of the aircraft, that is, with its mass, speed and height. This has a direct consequence for planning. A planner that minimizes mission time by increasing speed \cite{dong2017efficient} is, under this framework, also enlarging the ground-risk margin it is supposed to preserve. In other words, a purely geometric optimization criterion can produce an unintended regulatory violation.

To contain that risk, the SORA 2.5 methodology defines a set of pre-flight mitigation factors that lower the intrinsic ground risk of an operation and, with it, the probability of a person on the ground being struck, in some cases removing the need for a full certification process. Among these, the M1(C) mitigation---ground observation---is the one most directly tied to uninvolved people: it credits the capability of the crew, or of onboard systems, to observe most of the overflown area, detect uninvolved persons within it, and maneuver so that they are not overflown \cite{EASA_AMC_GM_SORA25_2025}.

 The present method does not implement the complete legal rules of either risk level. It implements the controller-level behavior that these rules motivate: when live perception suggests that the current local corridor should not be overflown, the UAV performs a temporary local bypass instead of continuing on the nominal line.

 
These regulatory concerns have already motivated a dedicated body of work. Primatesta \emph{et al.} combined off-line and on-line risk-aware planning driven by urban risk maps \cite{primatesta2019riskaware} and later systematized the construction of those ground-risk maps \cite{primatesta2020groundrisk}. Pang \emph{et al.} folded third-party risk into a metropolitan routing cost \cite{pang2022integrated}, while Tang \emph{et al.} modeled obstacle, fatality, and property-loss risk within a Min-cost A* search \cite{tang2023thirdparty}. Spatiotemporal and grid-based variants of ground-risk mapping have been explored by Pilko \emph{et al.} and Zhu \emph{et al.} \cite{pilko2023spatiotemporal,zhang2023groundrisk}. Closer to current European practice, Heinze and Uijt de Haag and, more recently, Bra{\ss}el \emph{et al.} encode SORA and target-level-of-safety notions directly into route optimization \cite{heinze2023sora,brassel2025riskaware}, and Du \emph{et al.} confirm that most civil-UAS safety research remains oriented toward assessment and authorization rather than runtime control \cite{du2024safetyrisk}.


%The most recent SOTA from 2024 to 2026 further sharpens the distinction between strategic risk planning and controller-level reaction. ElSayed and Mohamed (2024) integrate digital-twin airspace discretization with trajectory optimization under legal and operational constraints \cite{elsayed2024digitaltwin}. Pang \emph{et al.} (2024) model dynamic ground-risk uncertainty for safe drone delivery through stochastic route optimization \cite{pang2024stochastic}. Zhou \emph{et al.} (2025) propose a risk-based path-planning scheme for complex air--ground environments \cite{zhou2025riskbased}, while Riedel (2025) reviews detect-and-avoid standards and shows how strongly recent compliance discussions are still centered on air-risk, separation, and certification interfaces rather than overflight of uninvolved people on the ground \cite{riedel2025daa}. Checker \emph{et al.} (2025) review mission planning from a regulatory perspective, but at the swarm and mission-allocation level \cite{checker2025regulatory}. In 2026, Zhu \emph{et al.} move toward adaptive urban airspace allocation under uncertainty \cite{zhu2026adaptive}, and Nafees \emph{et al.} summarize the broader planning landscape while still highlighting real-time adaptation and constraint integration as open challenges \cite{nafees2026review}. Collectively, these works show that the field is advancing quickly, but the newest contributions still concentrate on strategic routing, digital twins, logistics airspace, or standards alignment more than on deterministic local replanning for inspection missions already underway.

\subsection{Research Gap}


%\begin{table*}[t]
%    \centering
%    \footnotesize
%    \begin{tabular}{@{}p{0.25\textwidth} p{0.32\textwidth} p{0.32\textwidth}@{}}
%        \toprule
%        \textbf{Representative work} & \textbf{What it already solves} & \textbf{Remaining gap relative to the proposed method} \\
%        \midrule
%        ElSayed and Mohamed (2024); Zhu \emph{et al.} (2026) \cite{elsayed2024digitaltwin,zhu2026adaptive}
%        & Digital-twin airspace structuring and adaptive urban allocation under legal, operational, and uncertainty constraints
%        & They operate mainly at city-airspace and logistics scale, not at the live local-replanning layer of an inspection UAV reacting to a detected civilian-like obstacle \\

%        Pang \emph{et al.} (2024); Zhou \emph{et al.} (2025); Bra{\ss}el \emph{et al.} (2025) \cite{pang2024stochastic,zhou2025riskbased,brassel2025riskaware}
%        & Strong third-party-risk, ground-risk, and TLS-aware optimization of routes
%        & Their main contribution is strategic or pre-flight route optimization rather than short-horizon mission-preserving replanning inside an active inspection controller \\

%        Riedel (2025); Checker \emph{et al.} (2025) \cite{riedel2025daa,checker2025regulatory}
%        & Up-to-date synthesis of detect-and-avoid standards and regulatory mission-planning constraints
%        & They clarify standards and regulatory context, but they do not provide a controller-level method for avoiding overflight of uninvolved people from live detections \\

%        Bellou \emph{et al.} (2024); Huang \emph{et al.} (2025) \cite{bellou2024aerial,huang2025hybrid}
%        & Inspection-specific perception and dynamic obstacle avoidance for power-line operations
%        & They improve inspection autonomy, but not around an explicit European no-overflight / ground-risk objective for uninvolved third parties \\

%        \textbf{This paper}
%        & \textbf{Deterministic waypoint-preserving local replanning triggered by live detections inside an autonomous inspection mission}
%        & \textbf{Brings protected-ground logic, controller integration, and auditability together in Simulated Environment rather than treating them as separate planning, perception, and compliance layers} \\
%        \bottomrule
%    \end{tabular}
%    \caption{Representative 2024--2026 comparison between recent SOTA strands and the proposed method.}
%    \label{tab:sota_positioning}
%\end{table*}

The proposed method is not the first UAV system to consider people on the ground, nor the first planner to use risk-aware costs. The contribution is more specific: existing inspection planners usually optimize geometric avoidance without an explicit uninvolved-people objective, whereas existing ground-risk and SORA-aware planners usually operate at strategic, route-design, or city-map scale. What remains weakly addressed is the controller-level problem of how an autonomous inspection UAS should react \emph{online}, during a running mission, when live detections suggest that the nominal corridor would overfly or closely pass above a civilian-like entity. The proposed method addresses exactly that intersection by: converting nearby detections into protected local no-fly regions, preserving the active mission waypoint as the anchor of the maneuver, and exposing the resulting decisions through the same telemetry. 


%To the best of our knowledge after reviewing the 2024--2026 literature, this makes the method a state-of-the-art contribution in the narrow niche of EASA-motivated, controller-level, inspection-integrated local replanning for uninvolved-people-aware UAS operation.

Among the available local-planning techniques, deterministic graph search is especially attractive for safety-critical systems because its completeness and failure modes are far easier to inspect than those of stochastic black-box policies. A* in particular offers a practical balance between computational simplicity and interpretable behavior \cite{hart1968formal}. The proposed method exploits that balance and combines it with a YOLO-based perception front-end \cite{ultralytics_yolo} and monocular obstacle geo-positioning inspired by the real-time approach of \cite{alaez2023real}.




%%%%%%%%%%%%%%%%%%%%%%%%%%%%%%%%%%%%%%%%%%

\section{Method and Materials}

This section is focused on the flow description of the proposed solution. As the model is designed for an UAS platform, it has to take into account every step involved in the flight, from the mission reading and takeoff of the aircraft, passing through the guided flight and conflict resolution, and finally the landing phase, once the mission is finished or caused by any problem encountered throughout the flight. 

The objective is to complete a powerline-inspection in an automatic way following a Flight Plan. But the special characteristic of this Flight Plan is that it is dynamic, and if any regulatory barrier or ordinary obstacle item (as a powerline asset) appears in the expected trajectory of the UAS, it is capable of performing a Flight Plan reconfiguration. 

The initial mission Flight Plan was generated based on the information available during the pre-flight phase, using the global planning tools and methodologies. In particular, the planning process relied on the nominal positions of the power-line assets, which may be outdated, incomplete, or inaccurately geo-referenced. The Flight Plan was defined using a map-based planning approach that incorporates additional constraints and contextual information, including communication network coverage, the availability of external services (e.g., GNSS or LTE), and previously known obstacle locations. These factors collectively influence the definition of Flight Plan waypoints (WPs),

Up to this point, all strategic mitigation measures that can be applied during the pre-flight phase have been addressed according to the corresponding operational risk assessment. These measures ensure compliance with the applicable European regulatory framework and consistency with the target operational category. 

\begin{figure}[h!]
\centering
\includegraphics[width=0.5\columnwidth]{Main/imagenes/System_Workflow.png}
\caption{System Workflow}
\label{solution}
\end{figure}

However, as introduced in Section 2, the operational focus shifts during the execution phase. At this stage, residual risks can arise due to dynamic, uncertain, or locally perceived conditions that cannot be fully captured during the planning phase. In these circumstances, simplified assumptions are often adopted and regulatory compliance is neglected. The proposed local reactive navigation model adapts the trajectory to the environment in real time.


The controller-level logic of the proposed method is specified below. At each iteration, the method receives the active waypoint, the latest perception outputs, and the current telemetry. Human-related detections are converted into protected local regions. If the nearest obstacle enters the dynamic reaction horizon, the method activates a bounded A* local bypass; otherwise, it preserves nominal mission following. 

\subsection{Problem Formulation}
At each control iteration, PORCE knows the current vehicle state, the active mission waypoint, and the most recent obstacle list. The method is only concerned with the local subproblem: move from the current position to the current waypoint while avoiding the obstacle subset that lies inside a bounded planning neighborhood.

Let the vehicle position be \(p_t \in \mathbb{R}^2\), the current waypoint be \(w_t \in \mathbb{R}^2\), and the perceived obstacles at time \(t\) be \(\mathcal{O}_t = \{o_1,\dots,o_k\}\). Each obstacle is inflated by a hard safety radius \(R_s\), generating a non-flyable region
\begin{equation}
    \mathcal{Z}_{\mathrm{safe}}(o_i) = \{q \in \mathbb{R}^2 : \lVert q - p_{o_i} \rVert_2 \le R_s\}.
\end{equation}

The local free space is therefore
\begin{equation}
    \mathcal{F}_t = \mathbb{R}^2 \setminus \bigcup_{i=1}^{k} \mathcal{Z}_{\mathrm{safe}}(o_i).
\end{equation}

In the current implementation, $(R_s = 12 m)$, which turns safety into a hard geometric constraint instead of a soft penalty.

Earlier conceptual project notes also discussed more general continuous cost-shaping ideas, including barrier-style formulations, as possible future directions. They are not the subject of the present paper. The system analyzed here is the one actually implemented in the repository: a deterministic bounded grid planner with hard obstacle inflation and explicit failsafe escalation, not a Control Barrier Function controller or any other unimplemented continuous optimizer.

\subsection{Dynamic Triggering Logic}

The method is not permanently active. The controller computes a dynamic reaction horizon
\begin{equation}
    D_{\mathrm{react}}(t) =
    \mathrm{clip}\!\left(D_{\mathrm{base}} + k_v \lVert v_t \rVert_2,\,
    D_{\min},\, D_{\max}\right),
\end{equation}
where $D_{\mathrm{base}}=45 meter$, $k_v= 2.0 s$, $D_{\min}=45 m$, and $D_{\max}=80 m$.

If the nearest perceived obstacle satisfies
\begin{equation}
    d_{\mathrm{nearest}}(t) \le D_{\mathrm{react}}(t),
\end{equation}
and the minimum replan interval has elapsed, PORCE requests a new local route from the replanning module. Otherwise, the vehicle continues nominal waypoint following.

\subsection{Bounded Local A* Planner}

The method uses a local occupancy grid centered on the vehicle. The validated configuration uses \(81 \times 81\) cells with $\delta=6 m$ resolution, yielding a bounded search window that is compatible with online execution inside the control loop.

Each cell \(c_{ij}\) is labeled as free or occupied according to the inflated obstacle set:
\begin{equation}
    S(c_{ij}) =
    \begin{cases}
        1, & \text{if } c_{ij} \cap \left(\mathbb{R}^2 \setminus \mathcal{F}_t\right) \neq \varnothing,\\
        0, & \text{otherwise.}
    \end{cases}
\end{equation}

The local route is computed with A* \cite{hart1968formal} using Euclidean heuristic,
\begin{equation}
    f(n) = g(n) + h(n),
\end{equation}
where \(g(n)\) is the cost of the accumulated path and \(h(n)\) is the Euclidean distance to the local goal. Eight-connectivity is enabled, so the method can combine cardinal and diagonal moves while remaining on the grid.

The target of the method is always the \emph{current mission waypoint}. This detail is important: the local planner is not allowed to redefine the global mission, only to create a temporary bypass that lets the UAV return to the original corridor as soon as the conflict disappears.

\subsection{Nominal Navigation, Temporary Inhibition, and Route Reattachment}

The method is best understood as a \emph{temporary inhibition} of nominal navigation, not as a new mission planner. 

Inside the control loop, the detection function is activated, so the system will continuously check the state of the obstacles with the detection function ($get\_detection$). The system responsible for providing the input to this function is an RGB camera on-boarded in the UAS. The methodology uses for detection of non regulatory compliant items or physical obstacles is based in the tracking proposed on \cite{alaez2023real}.

This detection function is what selects between the two possible paths inside the control loop. When no conflict is present, the system simply executes the nominal navigation flow along the previously planned trajectory, which works as follows: 
\begin{itemize}
    \item If $current\_idx$ is less than the length of the proposed flight plan, the system continue with the planned flight path. At each iteration, it checks whether the distance to the current target waypoint is below a threshold tolerance ($WP\_TOLERANCE$); if so, the target is updated to the next waypoint for the subsequent iteration. 
    \item If not, it implies that the last waypoint of the Flight Plan is reached, so the mission terminates normally through the landing logic.
\end{itemize}

Conversely, when the nearest obstacle crosses the dynamic reaction horizon, the nominal flow is suspended and an evasion path becomes the active short-horizon reference. During that interval, the aircraft no longer attempts to cut through the original corridor; instead, it follows the local sub-goals returned by the bounded A* planner. As soon as that local path is exhausted, nominal waypoint following resumes automatically toward the same mission waypoint that was active before the conflict.

This behavior is what makes the method operationally useful for power-line inspection. Its purpose is not merely to avoid collisions but to resolve local conflicts in real time and, once safe conditions are restored, to return to the globally planned route as quickly as possible. Operational safety is thus maintained without sacrificing the efficiency objective set during the initial global planning stage, whether that is optimized data acquisition, mission progression, or any other mission-specific metric established a priori.


\subsection{Failsafe Escalation}

If no valid local route is found, PORCE does not continue nominal flight blindly. Instead, the system integration includes an explicit escalation chain with hold, lateral replan, and terminal \texttt{LAND}/\texttt{RTL} actions. This failsafe design is part of the method even when it is not triggered in a particular experiment because it bounds the controller behavior under planner failure.

\subsection{Algorithmic Summary}

\begin{lstlisting}[language=Python, caption=Simplified control loop for the proposed method]
Flight_Plan = load_plan()
while Mission
    refresh_telemetry_and_obstacles()
    if telemetry_stale:
        continue

    if failsafe_action_active:
        execute_failsafe_action()
        continue

    wp = waypoints[current_wp_idx]
    reaction = clamp(base + groundspeed * gain,...
    min_r, max_r)
    nearest = nearest_obstacle()

    if nearest and nearest.dist < reaction and ...
    replan_allowed_now():
        route = astar_local(drone_pos, wp,... 
        planner_subset_obstacles())
        if valid(route):
            activate_evasion(route)
        else:
            apply_failsafe_escalation(nearest.dist)

    if evasion_active:
        follow_evasion_subpoints()
    else:
        follow_mission_waypoint(wp)
        if wp < length (Flight_Plan): 
            dist = get_WP_distance()
            if dist < WP_TOLERANCE_M
                current_wp_idx =+1
        else: 
            FlightMode = ['LAND']
            Mission = 0

\end{lstlisting}

\subsection{Method Integration in Simulation Workflow}

The proposed method is evaluated inside a validated \texttt{SIMULATION} workflow. %in the repository \cite{Deep_Aero_Twin}. 
The runtime architecture is specified in the figure \ref{Simulation Environment Diagram}. The mission is read from a QGroundControl waypoint file, vehicle state is received from SITL through MAVLink, obstacles are produced by the vision module, and the PORCE system coordinates nominal waypoint following, replanning activation, and visualization outputs.

\begin{figure}[h!]
\centering
\includegraphics[width=0.7\columnwidth]{Main/imagenes/Pipeline_A.png}
\caption{Simulation Environment Diagram}
\label{Simulation Environment Diagram}
\end{figure}

The Simulation Environment roles are the following:
\begin{itemize}
    \item \textbf{Mission input:} This is the file that contains the nominal inspection route.
    \item \textbf{SITL:} Simulation In The Loop tool.
    \item \textbf{MAVLINK Telemetry}:
    \item \textbf{Cesium / Unreal}: It is the visual Source of the environment. 
    \item \textbf{PORCE:} the PORCE system executes the control loop, manages waypoints, handles failsafe logic, and calls the replanning module whenever local replanning is necessary.
    \item \textbf{Perception:} This module runs a YOLO-based detector and projects detections in geographic coordinates using onboard state and the monocular geometry \cite{alaez2023real}. It is part of the PORCE control loop. 
    \item \textbf{Local Avoidance Planner:} This function implements the bounded A* local planner used to generate short-range evasion paths. It is part of the PORCE control loop. 
    \item \textbf{Visualization:} the visualization endpoint expose the evolution of the mission for analysis and debugging.
    \item \textbf{autopilot}: It runs inside SITL. Setpoint Execution. 
\end{itemize}


%\subsection{Runtime Integration with Simulation Workflow}
%The integration of the proposed method into Simulation Workflow is a closed-loop runtime contract rather than a loose module coupling. 
%\newline
%First, the mission is loaded from a QGroundControl waypoint file, and PORCE receives vehicle telemetry from SITL through MAVLink;
%\newline 
%Second, the perception subsystem queries the latest fused vehicle state and publishes detected obstacles to PORCE through \texttt{POST /api/obstacles}. PORCE then normalizes obstacle classes, enforces source filtering and token validation, maintains obstacle tracks with time-to-live logic, and rebuilds the active obstacle list used for decision making;
%\newline
%In the current implementation, human-related detections are already given a dedicated operational treatment. The labels \texttt{person}, \texttt{bicycle}, and \texttt{biker} are canonicalized into the same moving obstacle family, while \texttt{tower} and \texttt{cow} are treated as static classes for tracking and association purposes. This is not yet a full legal ontology of uninvolved people, but it is the concrete mechanism by which civilian-like detections become protected local regions that the method must avoid rather than overfly;
%\newline
%Third, the method is only evaluated when the autonomous mission is active, fresh obstacle evidence exists, and the nearest obstacle crosses a speed-adaptive reaction horizon. Before invoking the planner, PORCE forms a subset of bounded obstacles for local replanning: in the validated configuration, only obstacles inside \SI{55}{\meter} and up to 16 tracks are passed to the planner. This keeps the replanning problem short-range and explicitly tied to the current mission waypoint instead of allowing corridor-wide route redefinition; 

\begin{table}[h!]
    \centering
    \small
    \begin{tabular}{@{}p{0.3\columnwidth} p{0.2\columnwidth} p{0.4\columnwidth}@{}}
        \toprule
        \textbf{Parameter} & \textbf{Value} & \textbf{Role} \\
        \midrule
        Mission file & 12 waypoints & Nominal inspection route with takeoff and landing phases \\
        Control-loop period & 0.10 s & PORCE update period \\
        Cruise speed & 8 m/s & Nominal horizontal mission speed \\
        Grid cell size & 6 m & Local A* discretization step \\
        Grid radius & 40 cells & Local search radius, i.e. an $81 \times 81$ occupancy grid \\
        Obstacle safety radius & 12 m & Hard inflation radius around each obstacle \\
        Base reaction distance & 45 m & Minimum horizon for replanning activation \\
        Reaction gain & 2.0 s & Speed-dependent enlargement of reaction horizon \\
        Reaction range & 45--80 m & Clamped reaction distance interval \\
        Route-point tolerance & 3 m & Distance to accept a local replanning sub-goal as reached \\
        Failsafe threshold & 22 m & Distance used for hold/lateral replan escalation \\
        Failsafe hold & 2.5 s & Temporary hold action when replanning fails \\
        \bottomrule
    \end{tabular}
    \caption{Main configuration values of the validation environment and the proposed method used throughout the paper.}
    \label{tab:config}
\end{table}

%Fourth, PORCE calls a specific function, where the start is the current telemetry position and the goal is the current mission waypoint. The planner inflates each obstacle by the configured safety radius, discretizes the neighborhood into a local \(81 \times 81\) occupancy grid with \SI{6}{\meter} cells, runs bounded A*, and converts the resulting grid path back into geographic sub-points. If a valid route is found, PORCE stores it as the active evasion path; if not, the integration layer escalates through hold, lateral replan, and terminal failsafe actions;
%\newline

%Finally, the method outputs are consumed by the rest of the Simulated Environment through two channels. On the control side, PORCE sends each local replanning sub-point to the autopilot as MAVLink position-target commands until the evasion route is completed, at which point nominal waypoint following resumes automatically. On the observability side, the same state is exposed through the UI endpoint together with evasion metadata, planner parameters, tracked obstacles, and audit counters, so visualization tools, Unreal consumers, and zero-trust logs all observe the same decision chain. This shared contract is what makes the method not just a planner in isolation, but a reproducible system integrated into a full autonomous inspection pipeline.


Table~\ref{tab:config} summarizes the main runtime parameters used to evaluate the method inside the Simulation Workflow. These values come from the validated simulation defaults already present in the repository.

%%%%%%%%%%%%%%%%%%%%%%%%%%%%%%%%%%%%%%%%%%
\section{Results} 
This section presents the results obtaining during the validation of the proposed solution. 

%%%A AÑADIR EN EL PUNTO DE LOS RESULTADOS QUE SE CONSIDERE QUE APORTA MÁS. 


In a first step, the Mission is downloaded into the simulated UAS. The mission comes from a preplanned waypoints file that contains a takeoff waypoint, a sequence of predefined inspection waypoints and a landing command. For the experiment, the UAS climbs from $500 m$ MSL to approximately $523 m$ MSL, follows the inspection corridor, and lands automatically once the final waypoint is reached. This initial planned behaviour is represented in Figure \ref{Navigation Results} 1A, where it is appreciable that at in the initial stages, the UAS is navigating without any perturbation, following the predefined waypoints path. 


As explained before, while the UAS is in the air and, therefore, the proposed solution is running inside the embedded computer, the system is continuously analyzing the environment to detect if there is any obstacle or non regulatory compliant item that needs to be avoided. In Figure \ref{Navigation Results} 1B, it is represented the instant where the camera has already detected an obstacle (i.e. the obstacle is inside the Reaction Range), but the distance between the UAS and the obstacle is more than the Base Reaction Distance. Therefore, the system is not calculating yet an alternative evasion path, as the obstacle is only a potential hazard yet. The complete encounter, with the raw detection and tracking overlays produced by the vision front-end, is provided as Supplementary Video S1.

\begin{figure}[h!]
\centering
\includegraphics[width=\columnwidth]{Main/imagenes/fig_static_evasion_sequence.png}
\caption{Static obstacle evasion sequence observed in the simulation workflow. (\textbf{1A}) Nominal waypoint following along the inspection corridor. (\textbf{1B}) The perception front-end detects the static obstacles (cattle) while they remain outside the reaction distance; the planner stays in monitoring state. (\textbf{1C}) The nearest obstacle enters the speed-adaptive reaction horizon and the bounded local A* replanning is triggered. (\textbf{1D}) Active evasion maneuver: the vehicle follows the local bypass while both obstacles remain tracked. (\textbf{1E}) Bypass execution with conservative separation preserved. (\textbf{1F}) Conflict resolved: the vehicle rejoins the original corridor and resumes nominal mission.}
\label{Navigation Results}
\end{figure}
In this section, the results of the experiments performed with the environment defined in the simulation workflow are represented.

In contrast to the previous paragraph and figure, where the obstacle had already been detected but remained outside the safety radius, Figure \ref{Navigation Results} 1C represents a later stage of the encounter. At this instant, the obstacle has entered the reaction distance, triggering the trajectory replanning process. As a result, the system begins to compute an avoidance route intended to prevent a potential conflict and preserve the required separation margins.


The route proposed by the UAS is highlighted in orange, so it can be clearly distinguished from the original trajectory. Figure \ref{Navigation Results} 1D therefore shows the change from a monitoring state (the obstacle is detected but no action is required) to an active guidance state (the system evaluates the conflict and generates an alternative path).

In Figure \ref{Navigation Results} 1E, it is appreciable how the UA is performing the evasion maneuver through the alternative generated path, while the static identified element is remaining well detected and no conflict is reported. 


Figure \ref{Navigation Results} 1F shows the complete resolution of the conflict, once the UAS has successfully rejoined its original path and can continue with the data acquisition mission. At this stage, the avoidance maneuver has been completed and the vehicle has returned to nominal operation, resuming the planned inspection or monitoring task. This moment represents the final phase of the guidance response, where the temporary deviation introduced to maintain safe separation is no longer required and the mission can proceed under normal conditions.

\begin{figure}[h!]
\centering
\includegraphics[width=0.5\columnwidth]{Main/imagenes/transcurso_mision_nuevo_obs.png}
\caption{Middle Stages. Rejoining and new detections}
\label{Middle Stages}
\end{figure}

At this point, the incident is also recorded and reported by the system for subsequent analysis. Although this may result in partial data loss or gaps in the information collected from the asset during the maneuver, such limitations are considered acceptable within the operational framework. In safety-critical applications, preserving the integrity of the UAS and maintaining a safe distance from obstacles must take precedence over uninterrupted data collection. This figure therefore shows both the completion of the conflict-resolution process and the operational principle behind it: mission continuity is always subordinated to safety requirements.



Moreover, Figure \ref{Middle Stages} also shows that the process does not remain as an isolated event. A new detection can be observed after the UAS has returned to its nominal trajectory, demonstrating that the system is not limited to a single conflict-resolution cycle. The system can repeat this cycle (path following, detection, replanning, and recovery to the original route) as many times as the scenario requires. This shows that the guidance strategy can handle successive conflicts and resume the mission after each one.


\section{Conclusions}

This paper has presented PORCE, a deterministic real-time local replanning method for waypoint-preserving obstacle avoidance in autonomous power-line inspection missions, and has validated it end-to-end inside a georeferenced Unreal/Cesium simulation workflow. The central contribution is not obstacle avoidance itself, which is a mature field, but the integration of the European ground-risk logic into the online controller. Static obstacle detections are converted into protected local regions with hard safety radii, so the vehicle performs a conservative, auditable detour instead of overflying them. The global inspection mission is preserved and resumes automatically once the conflict is resolved.

The static-obstacle validation campaign shows three useful properties of the method. First, the planner remains inactive while no obstacle crosses the speed-adaptive reaction horizon, and it only intervenes when a conflict is confirmed, so the nominal mission is not perturbed unnecessarily. Second, the evasion is bounded and recoverable: the local A* bypass is computed over a small occupancy grid anchored to the current mission waypoint, produces a short conservative detour, and hands control back to the original corridor as soon as the protected region is cleared. Third, the whole decision chain, from the perception front-end to the evasion-route generation, can be reproduced from the repository artifacts. This traceability is a requirement in safety-critical aeronautical contexts.

The validation is based on a georeferenced digital twin, so every perception input, planner decision, and evasion trajectory can be traced and reproduced from the logged artifacts. This also allows a frame-level examination of the controller that is difficult to obtain in physical flight testing. Within this environment, the method applies the no-overflight principle through human-related obstacle classes and conservative geometric margins. In this way, part of the European ground-risk logic is encoded at the controller level, where certification-relevant online behavior must be guaranteed. The present campaign focuses on static obstacles, which directly matches the power-line asset and livestock scenarios most relevant to inspection operations; the treatment of moving actors is addressed in a companion study.

Future work builds directly on these results. The same deterministic methodology can be transferred to the real-flight workflow, where the absence of learned components simplifies certification. Ongoing work extends the scenario set with moving obstacles, defines class-dependent safety envelopes for people and other critical entities, and measures runtime on representative onboard hardware, with the goal of progressively extending the validated operational envelope.

%%%%%%%%%%%%%%%%%%%%%%%%%%%%%%%%%%%%%%%%%%

%\section{Conclusions}

%This section is not mandatory but can be added to the manuscript if the discussion is unusually long or complex.

%%%%%%%%%%%%%%%%%%%%%%%%%%%%%%%%%%%%%%%%%%

%%%%%%%%%%%%%%%%%%%%%%%%%%%%%%%%%%%%%%%%%%
\vspace{6pt} 

%%%%%%%%%%%%%%%%%%%%%%%%%%%%%%%%%%%%%%%%%%
%% optional
\supplementary{The following supporting information can be downloaded at \linksupplementary{s1}, Video S1: Annotated perception output of the cattle-evasion encounter shown in Figure \ref{Navigation Results}, recorded from the on-board camera of the simulated UAS with the detection and tracking overlays produced by the vision front-end (\texttt{vision\_yolo\_cow\_evasion\_v2\_FINAL.mp4}).}

% Only for journal Methods and Protocols:
% If you wish to submit a video article, please do so with any other supplementary material.
% \supplementary{The following supporting information can be downloaded at \linksupplementary{s1}, Figure S1: title; Table S1: title; Video S1: title. A supporting video article is available at doi: link.}

% Only used for preprtints:
% \supplementary{The following supporting information can be downloaded at the website of this paper posted on \href{https://www.preprints.org/}{Preprints.org}.}

% Only for journal Hardware:
% If you wish to submit a video article, please do so with any other supplementary material.
% \supplementary{The following supporting information can be downloaded at \linksupplementary{s1}, Figure S1: title; Table S1: title; Video S1: title.\vspace{6pt}\\
%\begin{tabularx}{\textwidth}{lll}
%\toprule
%\textbf{Name} & \textbf{Type} & \textbf{Description} \\
%\midrule
%S1 & Python script (.py) & Script of python source code used in XX \\
%S2 & Text (.txt) & Script of modelling code used to make Figure X \\
%S3 & Text (.txt) & Raw data from experiment X \\
%S4 & Video (.mp4) & Video demonstrating the hardware in use \\
%... & ... & ... \\
%\bottomrule
%\end{tabularx}
%}

%%%%%%%%%%%%%%%%%%%%%%%%%%%%%%%%%%%%%%%%%%


%\durcstatement{Current research is limited to the [please insert a specific academic field, e.g., XXX], which is beneficial [share benefits and/or primary use] and does not pose a threat to public health or national security. Authors acknowledge the dual-use potential of the research involving xxx and confirm that all necessary precautions have been taken to prevent potential misuse. As an ethical responsibility, authors strictly adhere to relevant national and international laws about DURC. Authors advocate for responsible deployment, ethical considerations, regulatory compliance, and transparent reporting to mitigate misuse risks and foster beneficial outcomes.}

% Only for journal Nursing Reports
%\publicinvolvement{Please describe how the public (patients, consumers, carers) were involved in the research. Consider reporting against the GRIPP2 (Guidance for Reporting Involvement of Patients and the Public) checklist. If the public were not involved in any aspect of the research add: ``No public involvement in any aspect of this research''.}
%
%% Only for journal Nursing Reports
%\guidelinesstandards{Please add a statement indicating which reporting guideline was used when drafting the report. For example, ``This manuscript was drafted against the XXX (the full name of reporting guidelines and citation) for XXX (type of research) research''. A complete list of reporting guidelines can be accessed via the equator network: \url{https://www.equator-network.org/}.}
%
%% Only for journal Nursing Reports
%\useofartificialintelligence{Please describe in detail any and all uses of artificial intelligence (AI) or AI-assisted tools used in the preparation of the manuscript. This may include, but is not limited to, language translation, language editing and grammar, or generating text. Alternatively, please state that “AI or AI-assisted tools were not used in drafting any aspect of this manuscript”.}




%%%%%%%%%%%%%%%%%%%%%%%%%%%%%%%%%%%%%%%%%%
%% Optional

%% Only for journal Encyclopedia
%\entrylink{The Link to this entry published on the encyclopedia platform.}

\authorcontributions{Conceptualization, P.D.P. and J.V.; methodology, P.D.P., X.O. and D.A.; software, X.O. and D.A.; validation, X.O. and D.A.; formal analysis, P.D.P. and J.V.; investigation, P.D.P. and X.O.; resources, P.D.P. and J.V.; data curation, X.O.; writing---original draft preparation, P.D.P. and X.O.; writing---review and editing, J.V. and D.A.; visualization, X.O.; supervision, J.V.; project administration, J.V. All authors have read and agreed to the published version of the manuscript.}

\funding{This research received no external funding.}

\institutionalreview{Not applicable.}

\informedconsent{Not applicable.}

\dataavailability{The data supporting the reported results are available from the corresponding author upon reasonable request.}

\conflictsofinterest{Pablo De Porcellinis is affiliated with FuVeX Civil SL. The remaining authors declare no conflicts of interest.}

\abbreviations{Abbreviations}{%
The following abbreviations are used in this manuscript:\\
\noindent
\begin{tabular}{@{}ll}
UAV & Unmanned Aerial Vehicle\\
UAS & Unmanned Aircraft System\\
PORCE & Path Planning and Obstacle Avoidance Regulatory Compliant Evasion\\
EASA & European Union Aviation Safety Agency\\
SORA & Specific Operations Risk Assessment\\
JARUS & Joint Authorities for Rulemaking on Unmanned Systems\\
GNSS & Global Navigation Satellite Systems\\
SITL & Software In The Loop\\
WP & Waypoint
\end{tabular}
}


\reftitle{References}
\bibliography{Main/referencias_revision_porce}
%} % If the paper is ``preprints'', please uncomment this parenthesis.
\end{document}

